\begin{figure}[htbp]
	\centering
	\begin{tikzpicture}[>=Stealth, line width=0.8pt, font=\small]
		
		\def\ang{30}    % angolo alpha^e in gradi
		\def\Lb{3.8}    % lunghezza asta
		\def\sep{6.0}   % offset pannello destro
		\def\ds{0.75}   % lunghezza frecce spostamento
		\def\fn{0.90}   % lunghezza frecce forza
		\def\fib{0.18}  % distanza fibra inferiore dall'asse
		
		%% ══════════════════════════════════════════════════════════════════
		%% PANNELLO (a) — CINEMATICA
		%% ══════════════════════════════════════════════════════════════════
		
		\coordinate (Pi) at (0, 0);
		\coordinate (Pj) at ({\Lb*cos(\ang)}, {\Lb*sin(\ang)});
		
		%% Fibra inferiore (tratteggiata): traslata di \fib nella dir. y_e = (sin,-cos)
		\coordinate (Pif) at ({\fib*sin(\ang)},              {-\fib*cos(\ang)});
		\coordinate (Pjf) at ({\Lb*cos(\ang)+\fib*sin(\ang)}, {\Lb*sin(\ang)-\fib*cos(\ang)});
		
		%% Assi globali
		\draw[->] (-1.0, -1.0) -- ++(1.3, 0) node[right] {$x$};
		\draw[->] (-1.0, -1.0) -- ++(0, 1.2) node[above] {$y$};
		
		%% Segmento orizzontale grigio (riferimento per alpha^e)
		\draw[gray!60, thin] (Pi) -- ++(1.8, 0);
		
		%% Arco alpha^e in grigio, raggio più grande e spostato a destra
		\draw[->, gray!60] (1.30, 0) arc (0:\ang:1.30);
		\node[gray!80] at ({1.60*cos(\ang/2)}, {0.20}) {$\alpha^e$};
		
		%% Fibra inferiore (tratteggiata, grigia)
		\draw[dashed, gray!70, line width=1.0pt] (Pif) -- (Pjf);
		
		%% Asta (linea d'asse continua)
		\draw[line width=1.4pt] (Pi) -- (Pj);
		
		%% Nodi (cerchietti vuoti)
		\filldraw[fill=white, draw=black, line width=1pt] (Pi) circle (3.5pt);
		\filldraw[fill=white, draw=black, line width=1pt] (Pj) circle (3.5pt);
		
		%% Etichette nodi
		\node[above left]  at (Pi) {$P_i$};
		\node[above left] at (Pj) {$P_j$};
		
		%% Spostamenti in P_i
		\draw[->, line width=1.1pt] (Pi) -- ++(\ds, 0)
		node[below] {$u_i$};
		\draw[->, line width=1.1pt] (Pi) -- ++(0, \ds)
		node[right]  {$v_i$};
		
		%% Spostamenti in P_j
		\draw[->, line width=1.1pt] (Pj) -- ++(\ds, 0)
		node[below] {$u_j$};
		\draw[->, line width=1.1pt] (Pj) -- ++(0, \ds)
		node[right] {$v_j$};
		
		%% Etichetta pannello
		\node[font=\small] at ({0.5*\Lb*cos(\ang)}, {-1.35}) {$(a)$};
		
		
		%% ══════════════════════════════════════════════════════════════════
		%% PANNELLO (b) — STATICA
		%% ══════════════════════════════════════════════════════════════════
		
		\coordinate (Pis) at (\sep, 0);
		\coordinate (Pjs) at ({\sep + \Lb*cos(\ang)}, {\Lb*sin(\ang)});
		
		%% Fibra inferiore pannello (b)
		\coordinate (Pisf) at ({\sep+\fib*sin(\ang)},                       {-\fib*cos(\ang)});
		\coordinate (Pjsf) at ({\sep+\Lb*cos(\ang)+\fib*sin(\ang)}, {\Lb*sin(\ang)-\fib*cos(\ang)});
		
		%% Assi globali
		\draw[->] ({\sep - 1.0}, -1.0) -- ++(1.3, 0) node[right] {$x$};
		\draw[->] ({\sep - 1.0}, -1.0) -- ++(0, 1.2) node[above] {$y$};
		
		%% Fibra inferiore (tratteggiata, grigia)
		\draw[dashed, gray!70, line width=1.0pt] (Pisf) -- (Pjsf);
		
		%% Asta
		\draw[line width=1.4pt] (Pis) -- (Pjs);
		
		%% Nodi (cerchietti vuoti)
		\filldraw[fill=white, draw=black, line width=1pt] (Pis) circle (3.5pt);
		\filldraw[fill=white, draw=black, line width=1pt] (Pjs) circle (3.5pt);
		
		%% Etichette nodi
		\node[above left]  at (Pis) {$P_i$};
		\node[above left] at (Pjs) {$P_j$};
		
		%% N^e in P_j: direzione +a_e (verso uscente estremo destro)
		\draw[->, line width=1.3pt] (Pjs)
		-- ++({cos(\ang)*\fn}, {sin(\ang)*\fn})
		node[below right] {$N^e$};
		
		%% N^e in P_i: direzione -a_e (verso uscente estremo sinistro)
		\draw[->, line width=1.3pt] (Pis)
		-- ++({-cos(\ang)*\fn}, {-sin(\ang)*\fn})
		node[below right] {$N^e$};
		
		%% Etichetta pannello
		\node[font=\small] at ({\sep + 0.5*\Lb*cos(\ang)}, {-1.35}) {$(b)$};
		
	\end{tikzpicture}
	\caption{Asta reticolare $e$ che connette i nodi $P_i$ (iniziale)
		e $P_j$ (finale), inclinata di $\alpha^e$ rispetto all'asse
		globale $x$. La linea tratteggiata indica la fibra inferiore
		e definisce il verso di percorrenza da $P_i$ a $P_j$.
		\emph{(a)}~Cinematica: spostamenti nodali $(u_i,\,v_i)$ e
		$(u_j,\,v_j)$ nel riferimento globale $(x,\,y)$.
		\emph{(b)}~Statica: sforzo normale $N^e$
		($N^e > 0$ indica trazione) alle due sezioni di
		estremità, con il verso convenzionale positivo rivolto
		verso l'esterno dell'asta.}
	\label{fig:asta_reticolare_cin_stat}
\end{figure}
